\documentclass[a4paper,11pt]{article}
\usepackage{lmodern}
\usepackage[T1]{fontenc}
\usepackage[utf8]{inputenc}
\usepackage{amsmath,amssymb,amsthm}
\usepackage{mathtools}
\usepackage{geometry}
\geometry{a4paper, margin=25mm}
\usepackage{hyperref}
\hypersetup{colorlinks=true, urlcolor=blue, linkcolor=blue, citecolor=blue}
\usepackage{enumitem}
\usepackage{array}
\usepackage{booktabs}
\usepackage{longtable}
\usepackage{graphicx}
\usepackage{fancyvrb}
\usepackage{xcolor}
\usepackage{calc}
\usepackage{etoolbox}

% Theorem environments (English)
\newtheorem{theorem}{Theorem}
\newtheorem{proposition}[theorem]{Proposition}
\newtheorem{lemma}[theorem]{Lemma}
\newtheorem{corollary}[theorem]{Corollary}
\theoremstyle{definition}
\newtheorem{definition}[theorem]{Definition}
\theoremstyle{remark}
\newtheorem{remark}[theorem]{Remark}
\newtheorem{observation}[theorem]{Observation}

\providecommand{\tightlist}{%
  \setlength{\itemsep}{0pt}\setlength{\parskip}{0pt}}

\providecommand{\pandocbounded}[1]{#1}
\setlength{\emergencystretch}{3em}

\title{Composition of Central Projection and the Closed Form of the
Composite Curvature Radius}
\date{An Algebraic Formulation of High-Dimensional Reduction via One
Central Projection and Commutative Cuts on the Sphere \\[1ex] 2026-05}
\author{Noriaki Kihara \\ \small WF System Co.,
Ltd. \\ \small ORCID: 0009-0004-6753-4020}

\begin{document}
\maketitle

\section{Abstract}\label{abstract}

Dimensional reduction by central projection from an \(n\)-dimensional
Euclidean space to a \(d\)-dimensional target space is often described
as ``the composition of successive central projections along multiple
axes.'' This paper points out that this description is generally
inaccurate and provides the correct formulation.

The main results are as follows:

\begin{itemize}
\tightlist
\item
  \textbf{First stage}: The central projection from the Euclidean
  background \(\mathbb{R}^n\) onto the \((n-1)\)-dimensional sphere
  \(S^{n-1}(r_1)\) of radius \(r_1\) occurs \textbf{only once}. This is
  the dimension-reducing map in the true sense.
\item
  \textbf{Second stage onward}: ``Axis-wise operations'' applied to a
  point already on the sphere \(S^{n-1}(r_1)\) are not central
  projections, but \textbf{simultaneity-section cuts on the sphere}.
  These operations only change the curvature radius and have no direct
  relation with the Euclidean background coordinates.
\item
  \textbf{Commutativity}: Cut operations along distinct axes on the
  sphere are completely \textbf{commutative}.
\item
  \textbf{Closed form of the composite curvature radius}: After
  successive cuts along the axis set
  \(S = \{i_1, \ldots, i_k\} \subseteq \{1, \ldots, n\}\), the final
  spherical radius is
\end{itemize}

\[
r_{\rm final}^{2} \;=\; r_{1}^{2} \;-\; \sum_{j=1}^{k}\bigl(x_{i_j}^{*}\bigr)^{2}
\]

where \(x_{i_j}^{*}\) is \textbf{the value of the axis \(x_{i_j}\)
component on the sphere immediately after the first central projection}.

\begin{itemize}
\tightlist
\item
  \textbf{Algebraic structure}: The set of axis-cut operations
  \(\{\sigma_i\}\) forms an Abelian semigroup, and the composite
  \(\sigma_S\) depends only on the axis subset \(S\).
\end{itemize}

This provides a unified description of dimensional reduction from any
\(n\)-dimensional space to a \(d < n\) dimensional space as a closed
operation on the combinatorics of axis selection.

Public information:

\begin{itemize}
\tightlist
\item
  DOI (Concept, auto-redirects to the latest version):
  \href{https://doi.org/10.5281/zenodo.20060728}{10.5281/zenodo.20060728}
\item
  DOI (v1):
  \href{https://doi.org/10.5281/zenodo.20060729}{10.5281/zenodo.20060729}
\item
  License: CC BY 4.0
\item
  Format: md / tex / pdf × JA/EN = 6 files
\end{itemize}

\begin{center}\rule{0.5\linewidth}{0.5pt}\end{center}

\section{1. Introduction}\label{introduction}

\subsection{1.1 Motivation}\label{motivation}

Central projection (also known as radial projection or gnomonic
projection) is a classical geometric operation that radially projects
points of a Euclidean space onto a sphere or a tangent plane {[}1{]}.

To realize dimensional reduction from an \(n\)-dimensional Euclidean
space to a \(d\)-dimensional target space (with \(d < n\)), one selects
\(k = n - d\) axes and applies central projection successively, as is
sometimes described {[}3{]}. For example, such descriptions take the
following form:

\begin{quote}
``Within an \(n\)-dimensional space, select axes
\(x_{i_1}, x_{i_2}, \ldots, x_{i_k}\) and apply central projection
successively along each axis direction to reduce to an
\((n-k)\)-dimensional space.''
\end{quote}

However, this description has an ambiguity that allows multiple
interpretations of the operation. Specifically, the fact that the first
operation and the second-and-subsequent operations are
\textbf{essentially different kinds of maps} is easily overlooked.

\subsection{1.2 Purpose of this paper}\label{purpose-of-this-paper}

The purpose of this paper is threefold:

\begin{enumerate}
\def\labelenumi{\arabic{enumi}.}
\tightlist
\item
  To clarify the essential difference between the first and the
  second-and-subsequent operations in the composition of central
  projections.
\item
  To show that the correct second-and-subsequent operation is ``a
  simultaneity-section cut on the sphere,'' and that this cut is
  commutative.
\item
  To derive the closed form of the composite curvature radius and
  establish the algebraic structure of the axis-cut operations (Abelian
  semigroup).
\end{enumerate}

\subsection{1.3 What this paper does not
address}\label{what-this-paper-does-not-address}

This paper is purely a paper on geometric algebraic structure. The
following are not addressed:

\begin{itemize}
\tightlist
\item
  The concrete interpretation of the axis \(x_i\) (its specific meaning
  in spacetime, observers, gauges, relativity, etc.)
\item
  The concrete meaning of the axis selection \(S\)
\item
  The physical necessity of specific dimensions \((n, d)\)
\item
  Generalization to non-Euclidean backgrounds (pseudo-Riemannian
  metrics, curved manifolds, etc.)
\end{itemize}

These are problems to be examined separately, taking the algebraic
results of this paper as given.

\begin{center}\rule{0.5\linewidth}{0.5pt}\end{center}

\section{2. Formulation of Central Projection and Cut
Operations}\label{formulation-of-central-projection-and-cut-operations}

\subsection{\texorpdfstring{2.1 Central projection \(\pi\): from
Euclidean background to
sphere}{2.1 Central projection \textbackslash pi: from Euclidean background to sphere}}\label{central-projection-pi-from-euclidean-background-to-sphere}

We define the \textbf{central projection} as the map that sends a point
\(P = (x_1, \ldots, x_n)\) of \(n\)-dimensional Euclidean space
\(\mathbb{R}^n\) (with \(P \neq 0\)) radially to the
\((n-1)\)-dimensional sphere of radius \(r_1\) centered at the origin
\(O\):

\[
S^{n-1}(r_1) \;:=\; \bigl\{X \in \mathbb{R}^{n} \;\big|\; \lVert X\rVert = r_1\bigr\}
\]

via the formula:

\[
\pi(P) \;=\; \frac{r_1}{\lVert P\rVert} P
\]

This is a well-defined map from \(\mathbb{R}^n \setminus \{O\}\) to
\(S^{n-1}(r_1)\).

The central projection \(\pi\) has the following properties:

\begin{itemize}
\tightlist
\item
  \textbf{(P1) Constraint on the image}: \(\lVert \pi(P) \rVert = r_1\)
  (for any \(P\)).
\item
  \textbf{(P2) Direction preservation}: \(\pi(P)\) has the same radial
  direction as \(P\).
\item
  \textbf{(P3) Dimension of the image}: The image \(\pi(P)\) is a point
  on the \((n-1)\)-dimensional manifold \(S^{n-1}(r_1)\).
\end{itemize}

\subsection{2.2 Axis components of points on the
sphere}\label{axis-components-of-points-on-the-sphere}

The point after central projection \(P' = \pi(P) \in S^{n-1}(r_1)\) has
\(n\) coordinate components in \(\mathbb{R}^n\):

\[
P' \;=\; (x_1', x_2', \ldots, x_n'), \qquad x_i' \;=\; \frac{r_1}{\lVert P\rVert} x_i
\]

We call the component \(x_i'\) corresponding to the axis \(x_i\) the
\textbf{axis \(x_i\) component on the sphere}, and denote it \(x_i^*\)
in what follows:

\[
x_i^{*} \;:=\; \frac{r_1}{\lVert P\rVert} x_i
\]

This quantity is uniquely determined once \(P\) is fixed.

\subsection{\texorpdfstring{2.3 Simultaneity-section cuts \(\sigma_i\)
on the
sphere}{2.3 Simultaneity-section cuts \textbackslash sigma\_i on the sphere}}\label{simultaneity-section-cuts-sigma_i-on-the-sphere}

Consider the operation of fixing the axis \(x_i\) component of a point
\(X\) on the sphere \(S^{n-1}(r)\) to \(c\). That is, take the
intersection of the hyperplane
\(\Sigma_i(c) := \{X \in \mathbb{R}^n \mid x_i = c\}\) with the sphere
\(S^{n-1}(r)\):

\[
S^{n-1}(r) \cap \Sigma_i(c) \;=\; \bigl\{X \;\big|\; \lVert X\rVert = r,\; x_i = c\bigr\}
\]

When \(|c| \leq r\), this intersection is isomorphic to the
\((n-2)\)-dimensional sphere of radius

\[
r' \;=\; \sqrt{r^{2} - c^{2}}.
\]

We identify this with the sphere \(S^{n-2}(r')\) within the
\((n-1)\)-dimensional space obtained by removing the axis \(x_i\).

We call this operation the \textbf{cut along the axis \(x_i\)} and
denote it \(\sigma_i\):

\[
\sigma_{i}: S^{n-1}(r) \;\longrightarrow\; S^{n-2}(r'), \qquad r' = \sqrt{r^{2} - c^{2}}
\]

Here \(c\) is a fixed value, and in this paper we adopt
\textbf{\(x_i^{*}\), the axis \(x_i\) component on the sphere
immediately after the first central projection}, as \(c\).

\subsection{2.4 Essential difference between central projection and
cut}\label{essential-difference-between-central-projection-and-cut}

{\def\LTcaptype{none} % do not increment counter
\begin{longtable}[]{@{}
  >{\raggedright\arraybackslash}p{(\linewidth - 4\tabcolsep) * \real{0.3333}}
  >{\raggedright\arraybackslash}p{(\linewidth - 4\tabcolsep) * \real{0.3333}}
  >{\raggedright\arraybackslash}p{(\linewidth - 4\tabcolsep) * \real{0.3333}}@{}}
\toprule\noalign{}
\begin{minipage}[b]{\linewidth}\raggedright
Property
\end{minipage} & \begin{minipage}[b]{\linewidth}\raggedright
Central projection \(\pi\)
\end{minipage} & \begin{minipage}[b]{\linewidth}\raggedright
Cut \(\sigma_i\)
\end{minipage} \\
\midrule\noalign{}
\endhead
\bottomrule\noalign{}
\endlastfoot
Input & Point of Euclidean background \(\mathbb{R}^n\) & Point on the
sphere \(S^{n-1}(r)\) \\
Output & Point on the sphere \(S^{n-1}(r_1)\) & Point on the sphere
\(S^{n-2}(r')\) \\
Dimension reduction & \(n \to (n-1)\) manifold & \((n-1) \to (n-2)\)
manifold \\
Radius & \(r_1\) (given) & \(r' = \sqrt{r^2 - c^2}\) (determined) \\
Reference to Euclidean background & Direct & None (intrinsic to the
sphere) \\
\end{longtable}
}

A central projection is ``a map from the Euclidean background to a
curved surface,'' whereas a cut is ``a geometric operation on a curved
surface.'' They are sometimes both referred to as ``dimension-reduction
operations,'' but they are \textbf{essentially different kinds of maps}.

\begin{center}\rule{0.5\linewidth}{0.5pt}\end{center}

\section{3. Commutativity of Cut
Operations}\label{commutativity-of-cut-operations}

\subsection{3.1 Main theorem:
commutativity}\label{main-theorem-commutativity}

\textbf{Theorem 1 (Commutativity of cuts).} For two distinct axes
\(x_i, x_j\) (\(i \neq j\)) on the sphere \(S^{n-1}(r)\), the cut
operations \(\sigma_i, \sigma_j\) commute:

\[
\sigma_j \circ \sigma_i \;=\; \sigma_i \circ \sigma_j
\]

\textbf{Proof.} The action of \(\sigma_i\) on a point
\(X = (x_1, \ldots, x_n) \in S^{n-1}(r)\) is:

\begin{itemize}
\tightlist
\item
  Fix the axis \(x_i\) component to \(c_i\) (with \(c_i = x_i^{*}\)).
\item
  The remaining point in the \((n-1)\)-dimensional space is \(X\) with
  the axis \(x_i\) removed.
\item
  The remaining point lies on the sphere \(S^{n-2}(r')\) where
  \(r' = \sqrt{r^2 - c_i^2}\).
\end{itemize}

The axis \(x_j\) component \(c_j\) \textbf{does not change} under the
cut \(\sigma_i\) (since the operation removes only the axis \(x_i\), the
values along other axes are preserved).

Order 1: \(\sigma_j \circ \sigma_i\):

\begin{itemize}
\tightlist
\item
  Cut \(\sigma_i\) changes the radius from \(r\) to
  \(\sqrt{r^2 - c_i^2}\) and removes the axis \(x_i\).
\item
  Cut \(\sigma_j\) changes the radius from \(\sqrt{r^2 - c_i^2}\) to
  \(\sqrt{r^2 - c_i^2 - c_j^2}\) and removes the axis \(x_j\).
\item
  Final radius: \(\sqrt{r^2 - c_i^2 - c_j^2}\).
\end{itemize}

Order 2: \(\sigma_i \circ \sigma_j\):

\begin{itemize}
\tightlist
\item
  Cut \(\sigma_j\) changes the radius from \(r\) to
  \(\sqrt{r^2 - c_j^2}\) and removes the axis \(x_j\).
\item
  Cut \(\sigma_i\) changes the radius from \(\sqrt{r^2 - c_j^2}\) to
  \(\sqrt{r^2 - c_j^2 - c_i^2}\) and removes the axis \(x_i\).
\item
  Final radius:
  \(\sqrt{r^2 - c_j^2 - c_i^2} = \sqrt{r^2 - c_i^2 - c_j^2}\).
\end{itemize}

Both yield the same radius, and the remaining coordinate components are
the same (the residue after removing both axes \(x_i\) and \(x_j\)).
Hence the maps coincide. \(\Box\)

\subsection{\texorpdfstring{3.2 Generalization: commutativity of
\(k\)-axis
cuts}{3.2 Generalization: commutativity of k-axis cuts}}\label{generalization-commutativity-of-k-axis-cuts}

\textbf{Corollary 1 (Commutativity of \(k\)-axis cuts).} For \(k\)
pairwise distinct axes \(x_{i_1}, \ldots, x_{i_k}\), the cuts
\(\sigma_{i_1}, \ldots, \sigma_{i_k}\) yield the same map under any
permutation:

\[
\sigma_{i_{\tau(1)}} \circ \sigma_{i_{\tau(2)}} \circ \cdots \circ \sigma_{i_{\tau(k)}} \;=\; \sigma_{i_1} \circ \sigma_{i_2} \circ \cdots \circ \sigma_{i_k}
\]

where \(\tau\) is any permutation of \(\{1, \ldots, k\}\).

\textbf{Proof.} By Theorem 1, any adjacent pair commutes, and the
symmetric group \(S_k\) is generated by adjacent transpositions.
\(\Box\)

\subsection{3.3 Composite cut by a set of
axes}\label{composite-cut-by-a-set-of-axes}

\textbf{Definition 2 (Composite cut).} For an axis set
\(S = \{i_1, \ldots, i_k\} \subseteq \{1, \ldots, n\}\), define the
composite cut \(\sigma_S\) by:

\[
\sigma_S \;:=\; \sigma_{i_1} \circ \sigma_{i_2} \circ \cdots \circ \sigma_{i_k}
\]

By Corollary 1, this is independent of the order and is determined
solely by \(S\).

\begin{center}\rule{0.5\linewidth}{0.5pt}\end{center}

\section{4. Closed Form of the Composite Curvature
Radius}\label{closed-form-of-the-composite-curvature-radius}

\subsection{4.1 Main theorem: the composite curvature
radius}\label{main-theorem-the-composite-curvature-radius}

\textbf{Theorem 2 (Composite curvature radius).} After being projected
onto the sphere \(S^{n-1}(r_1)\) of radius \(r_1\) by the first central
projection, applying the composite cut \(\sigma_S\) for the axis set
\(S = \{i_1, \ldots, i_k\}\) yields the final sphere of radius:

\[
\boxed{\;r_{\rm final}^{2} \;=\; r_{1}^{2} \;-\; \sum_{j=1}^{k}\bigl(x_{i_j}^{*}\bigr)^{2}\;}
\]

where \(x_{i_j}^{*}\) is the axis \(x_{i_j}\) component on the sphere
immediately after the first central projection.

\textbf{Proof.} In the proof of Theorem 1 (or Corollary 1), each cut
decreases \(r^2\) by \((x_{i_j}^{*})^2\). Repeating this \(k\) times
yields:

\[
r_{\rm final}^{2} \;=\; r_{1}^{2} - (x_{i_1}^{*})^{2} - (x_{i_2}^{*})^{2} - \cdots - (x_{i_k}^{*})^{2}
\]

Since the result is independent of order, the closed form holds.
\(\Box\)

\subsection{4.2 Pythagorean structure}\label{pythagorean-structure}

The closed form for the composite curvature radius takes the form of
\textbf{a sum of \((x_{i_j}^{*})^2\)} over the cut axes. This is nothing
other than the higher-dimensional version of the Pythagorean theorem.

Specifically, for a point \(P = (x_1, \ldots, x_n)\) on the sphere
\(S^{n-1}(r_1)\) centered at the origin,

\[
r_1^{2} \;=\; \sum_{i=1}^{n} x_i^{2}
\]

When we apply cuts along the axis set \(S\), the axis \(x_{i_j}\)
components \((x_{i_j}^{*})\) are ``subtracted'' from the squared sphere
radius, and the remaining \((n-k)\) axis components form the new
sphere's squared radius:

\[
r_{\rm final}^{2} \;=\; \sum_{i \notin S} x_i^{2} \;=\; r_1^{2} - \sum_{i \in S} (x_i^{*})^{2}
\]

This means that ``the coordinate components of a point on the sphere
contribute independently as an orthogonal decomposition by the axis
set.''

\subsection{4.3 Invariance of remaining
coordinates}\label{invariance-of-remaining-coordinates}

\textbf{Corollary 2 (Invariance of remaining coordinates).} Even after
applying the composite cut \(\sigma_S\) for the axis set \(S\), the
coordinate components of axes not in \(S\) do not change:

\[
\bigl(\sigma_S(P)\bigr)_i \;=\; x_i \qquad \text{for all } i \notin S
\]

\textbf{Proof.} Each cut \(\sigma_i\) is the operation that removes the
axis \(x_i\) component, leaving the components of other axes untouched.
Hence, even after applying the entire composite \(\sigma_S\), the
components of axes not in \(S\) remain invariant. \(\Box\)

\subsection{4.4 Invertibility of the composite
cut}\label{invertibility-of-the-composite-cut}

\textbf{Corollary 3 (Invertibility).} The composite cut \(\sigma_S\)
admits an inverse map: from the set of cut-axis values
\(\{x_i^{*}\}_{i \in S}\) and the image point \(\sigma_S(P)\) on the
post-cut sphere \(S^{n-|S|-1}(r_{\rm final})\), the original point \(P\)
on the pre-cut sphere \(S^{n-1}(r_1)\) is uniquely recovered.

\textbf{Proof.} The image point \(\sigma_S(P)\) has \((n - |S|)\)
coordinate components, each unchanged from before the cut by Corollary
2. The cut-axis values \(x_i^{*}\) (\(i \in S\)) are also preserved.
Hence all \(n\) coordinate components of the original point \(P\) on
\(S^{n-1}(r_1)\) are recoverable. The radius is uniquely determined by
Theorem 2 as \(r_1^2 = r_{\rm final}^2 + \sum (x_i^{*})^2\). \(\Box\)

\textbf{Note}: The corollary holds even in the degenerate case
\(k = n - 1\) (where the final sphere is 0-dimensional, i.e., two points
on a 1-dimensional number line).

\begin{center}\rule{0.5\linewidth}{0.5pt}\end{center}

\section{5. Algebraic Structure: Abelian
Semigroup}\label{algebraic-structure-abelian-semigroup}

\subsection{5.1 Composition by set
operations}\label{composition-by-set-operations}

\textbf{Theorem 3 (Composition by set operations).} For axis sets
\(S, T \subseteq \{1, \ldots, n\}\) with \(S \cap T = \emptyset\):

\[
\sigma_{S \cup T} \;=\; \sigma_S \circ \sigma_T
\]

\textbf{Proof.} By definition, both sides are the composition of cuts
along all axes in \(S \cup T\). By Corollary 1, the result is
independent of order, so both sides are equal. \(\Box\)

\subsection{5.2 Semigroup structure}\label{semigroup-structure}

The set of axis-cut operations
\(\{\sigma_S \mid S \subseteq \{1, \ldots, n\}\}\) has the following
structure:

\begin{itemize}
\tightlist
\item
  \textbf{Operation}: composition \(\circ\)
\item
  \textbf{Commutativity}:
  \(\sigma_S \circ \sigma_T = \sigma_T \circ \sigma_S\) when
  \(S \cap T = \emptyset\)
\item
  \textbf{Associativity}:
  \((\sigma_S \circ \sigma_T) \circ \sigma_U = \sigma_S \circ (\sigma_T \circ \sigma_U)\)
  for pairwise disjoint \(S, T, U\)
\item
  \textbf{Identity}: \(\sigma_\emptyset = \mathrm{id}\) (identity map)
\end{itemize}

This is the structure of an \textbf{Abelian semigroup} {[}2{]}
isomorphic to subset union (when disjoint). Inverse elements do not
exist (an inverse operation that increases dimension is not defined), so
this is a semigroup, not a group.

\subsection{5.3 Integration of the overall
structure}\label{integration-of-the-overall-structure}

The total reduction from an \(n\)-dimensional Euclidean space to a
\(d\)-dimensional final space is:

\begin{enumerate}
\def\labelenumi{\arabic{enumi}.}
\tightlist
\item
  \textbf{First stage (central projection)}:
  \(\pi: \mathbb{R}^n \to S^{n-1}(r_1)\), exactly once.
\item
  \textbf{Second stage (composite cut)}:
  \(\sigma_S: S^{n-1}(r_1) \to S^{d-1}(r_{\rm final})\), with
  \(|S| = n - d\) axes.
\end{enumerate}

The composite total map is:

\[
\Phi_{S} \;=\; \sigma_S \circ \pi \;:\; \mathbb{R}^{n} \setminus \{O\} \;\longrightarrow\; S^{d-1}(r_{\rm final})
\]

The final radius is computable in closed form as
\(r_{\rm final}^2 = r_1^2 - \sum_{i \in S}(x_i^{*})^2\).

\begin{center}\rule{0.5\linewidth}{0.5pt}\end{center}

\section{6. Numerical Verification}\label{numerical-verification}

\subsection{\texorpdfstring{6.1 Example: 7-dimensional \(\to\)
5-dimensional}{6.1 Example: 7-dimensional \textbackslash to 5-dimensional}}\label{example-7-dimensional-to-5-dimensional}

For the 7-dimensional Euclidean point \(P = (1, 2, 3, 4, 10, 5, 3)\), we
apply central projection with \(r_1 = 10\), then a composite cut along 2
axes (\(x_5, x_6\)).

\textbf{Central projection}:

\(\lVert P \rVert = \sqrt{1 + 4 + 9 + 16 + 100 + 25 + 9} = \sqrt{164} \approx 12.806\)

After central projection:
\(P' = (10/\sqrt{164}) \cdot P \approx (0.781, 1.562, 2.343, 3.123, 7.809, 3.904, 2.343)\)

Axis \(x_5\) component: \(x_5^{*} \approx 3.904\) (the original value 5
multiplied by \(r_1 / \lVert P \rVert\))

Axis \(x_6\) component: \(x_6^{*} \approx 2.343\) (the original value 3
multiplied by \(r_1 / \lVert P \rVert\))

\textbf{Composite cut}:

Order A (\(x_5 \to x_6\)):

\begin{itemize}
\tightlist
\item
  Radius after the \(x_5\) cut:
  \(\sqrt{10^2 - 3.904^2} = \sqrt{100 - 15.244} \approx \sqrt{84.756} \approx 9.206\)
\item
  Radius after the \(x_6\) cut:
  \(\sqrt{9.206^2 - 2.343^2} = \sqrt{84.756 - 5.488} \approx \sqrt{79.268} \approx 8.903\)
\end{itemize}

Order B (\(x_6 \to x_5\)):

\begin{itemize}
\tightlist
\item
  Radius after the \(x_6\) cut:
  \(\sqrt{10^2 - 2.343^2} = \sqrt{100 - 5.488} \approx \sqrt{94.512} \approx 9.722\)
\item
  Radius after the \(x_5\) cut:
  \(\sqrt{9.722^2 - 3.904^2} = \sqrt{94.512 - 15.244} \approx \sqrt{79.268} \approx 8.903\)
\end{itemize}

Both orders yield the same final radius \(r_{\rm final} \approx 8.903\).

\textbf{Verification by closed form}:

\[
r_{\rm final}^{2} \;=\; r_{1}^{2} - (x_5^{*})^{2} - (x_6^{*})^{2} \;=\; 100 - 15.244 - 5.488 \;=\; 79.268
\]

\[
r_{\rm final} \;=\; \sqrt{79.268} \;\approx\; 8.903
\]

Order A, Order B, and the closed form all agree.

\subsection{6.2 Verification of invariance of remaining
coordinates}\label{verification-of-invariance-of-remaining-coordinates}

The 5-dimensional components of the final point after the composite cut:

\[
(0.781, 1.562, 2.343, 3.123, 7.809)
\]

These are the first 5 components of \(P'\) immediately after central
projection, unchanged by the cut operations. This verifies Corollary 2.

\subsection{6.3 Verification of
invertibility}\label{verification-of-invertibility}

Using the 5-dimensional components
\((0.781, 1.562, 2.343, 3.123, 7.809)\) and the cut-axis values
\(x_5^{*} = 3.904\), \(x_6^{*} = 2.343\), we recover the point on the
pre-cut sphere \(S^6(r_1 = 10)\):

\[
P' \;=\; (0.781, 1.562, 2.343, 3.123, 7.809, 3.904, 2.343)
\]

This is the original \(P'\), and \(\lVert P' \rVert = 10 = r_1\) is also
recovered. This verifies Corollary 3.

\begin{center}\rule{0.5\linewidth}{0.5pt}\end{center}

\section{7. Compatibility with Existing
Descriptions}\label{compatibility-with-existing-descriptions}

\subsection{7.1 Interpretation of ``successive central projections along
multiple
axes''}\label{interpretation-of-successive-central-projections-along-multiple-axes}

When the reduction from \(n\)-dimensional to \(d\)-dimensional is
described as ``successive central projections along multiple axes,'' the
formulation in this paper interprets it as follows:

\begin{itemize}
\tightlist
\item
  ``First central projection'': The actual central projection \(\pi\)
  from the Euclidean background to the sphere.
\item
  ``Second-and-subsequent central projections'': Cut operations
  \(\sigma_i\) on the sphere.
\end{itemize}

It is possible to refer to both as ``central projection'' without
distinction, but in that case:

\begin{itemize}
\tightlist
\item
  First: \(\mathbb{R}^n \to S^{n-1}(r_1)\), true dimension reduction.
\item
  Second-and-subsequent: \(S^m(r) \to S^{m-1}(r')\), dimension reduction
  between spheres.
\end{itemize}

One should note that this two-stage structure is implicitly assumed.

\subsection{7.2 Note on order
dependence}\label{note-on-order-dependence}

If ``projection along axis \(x_i\)'' in the Euclidean background is
naively interpreted as ``dividing all coordinates by \(x_i\),'' then
composition along multiple axes can become order-dependent.
Specifically, treating a point on the sphere (after projection along
axis \(x_i\)) again as a point of the Euclidean background and trying to
project along axis \(x_j\) leads to ambiguity in the meaning of the
operation.

In the framework of this paper, after the first central projection,
points are treated as ``points on the sphere,'' and
second-and-subsequent operations apply the cut \(\sigma_j\) on the
sphere. With this distinction, composition becomes completely
commutative.

\begin{center}\rule{0.5\linewidth}{0.5pt}\end{center}

\section{8. Conclusion}\label{conclusion}

This paper has established:

\begin{itemize}
\tightlist
\item
  Dimensional reduction by central projection essentially consists of
  \textbf{one true central projection plus \(k\) cut operations on the
  sphere}.
\item
  The cut operations on the sphere are completely \textbf{commutative},
  and a closed composition operation \(\sigma_S\) for axis sets is
  defined.
\item
  The composite curvature radius is given in closed form
\end{itemize}

\[
r_{\rm final}^{2} \;=\; r_{1}^{2} - \sum_{i \in S}(x_i^{*})^{2}
\]

with a Pythagorean orthogonal decomposition structure.

\begin{itemize}
\tightlist
\item
  The set of cut operations forms an Abelian semigroup.
\item
  Remaining coordinates are invariant under cut operations.
\item
  The composite cut is \textbf{invertible}: from the cut-axis values and
  the image point, the original point on the sphere can be uniquely
  recovered.
\end{itemize}

These algebraic structures function as a foundational language in
various geometric and physical frameworks dealing with reduction from
high to low dimensions.

\begin{center}\rule{0.5\linewidth}{0.5pt}\end{center}

\section{9. Open Problems}\label{open-problems}

This paper is purely a paper on geometric algebraic structure. The
following problems lie outside its scope and are to be examined
separately:

\begin{itemize}
\tightlist
\item
  \textbf{Concrete interpretation}: The specific meaning of axes
  \(x_i\), the necessity of particular \((n, d)\), and relations to
  observers.
\item
  \textbf{Algebraic structure of the inverse operation}: Invertibility
  (Corollary 3) exists at the level of maps; whether to extend this into
  a group structure on the algebraic structure of cut operations, or
  remain at the semigroup level, is a matter of choice.
\item
  \textbf{Non-Euclidean backgrounds}: Analogous structures for
  pseudo-Riemannian metrics or curved background manifolds.
\item
  \textbf{Group-theoretic extensions}: Group-theoretic structures of
  transformations that change the axis selection (continuous choices of
  axes), and correspondences with classical transformation groups.
\item
  \textbf{Continuous cut values}: This paper fixes the cut value \(c\)
  to the spherical component \(x_i^{*}\) immediately after the first
  central projection; an algebraic structure for cuts with more general
  \(c\) is left open.
\end{itemize}

\begin{center}\rule{0.5\linewidth}{0.5pt}\end{center}

\section{References}\label{references}

{[}1{]} Snyder, J.P. (1987). \emph{Map Projections---A Working Manual}.
U.S. Geological Survey Professional Paper 1395, U.S. Government Printing
Office, Washington, D.C. (Modern standard reference for map projections
including gnomonic projection. Cited in §1.1, §2.1 as the source of the
classical definition of central projection.)

{[}2{]} Howie, J.M. (1995). \emph{Fundamentals of Semigroup Theory}.
London Mathematical Society Monographs, New Series, Vol. 12, Oxford
University Press. (Standard reference for semigroup theory. Cited in §5
as the source of the Abelian semigroup structure.)

{[}3{]} Kihara, N. (2026). \emph{A Geometric Formulation of
4-Dimensional Space via Central Projection} (in Japanese). Zenodo. DOI:
\href{https://doi.org/10.5281/zenodo.19427780}{10.5281/zenodo.19427780}.
(Foundational paper of the central projection framework. Cited in §1.1
as one of the sources of the informal ``successive central projections
along multiple axes'' description that this paper aims to correct. The
algebraic results of the present paper hold independently of the
application contexts in the cited paper.)

\begin{center}\rule{0.5\linewidth}{0.5pt}\end{center}

\section{Author Information}\label{author-information}

\begin{itemize}
\tightlist
\item
  \textbf{Author}: Noriaki Kihara
\item
  \textbf{Affiliation}: WF System Co., Ltd.~/ Graduated from the School
  of Engineering Science, Osaka University (1983)
\item
  \textbf{ORCID}:
  \href{https://orcid.org/0009-0004-6753-4020}{0009-0004-6753-4020}
\item
  \textbf{License}: CC BY 4.0
\end{itemize}

\begin{center}\rule{0.5\linewidth}{0.5pt}\end{center}

\textbf{Note}: This paper provides an independent algebraic foundation
for dimensional reduction by central projection. Its algebraic results
do not depend on specific dimensions, axis selections, or physical
interpretations, and provide a universal foundation for any application
that requires them. Specific applications are developed in separate
papers.

\end{document}
